% Generated draft: requires mathematical and source review.
\documentclass[11pt]{article}
\usepackage[margin=25mm]{geometry}
\usepackage{fontspec}
\setmainfont{Latin Modern Roman}
\usepackage{amsmath,amssymb,mathtools,unicode-math}
\setmathfont{Latin Modern Math}
\usepackage{enumitem,needspace,xurl,hyperref}
\setlength{\parindent}{0pt}
\setlength{\parskip}{5pt}
\setlength{\emergencystretch}{3em}
\newcommand{\N}{\mathbb{N}}
\newcommand{\Z}{\mathbb{Z}}
\newcommand{\Q}{\mathbb{Q}}
\newcommand{\R}{\mathbb{R}}
\newcommand{\C}{\mathbb{C}}
\newcommand{\F}{\mathbb{F}}
\newcommand{\A}{\mathbb{A}}
\newcommand{\PP}{\mathbb{P}}
\newcommand{\E}{\mathbb{E}}
\newcommand{\eps}{\varepsilon}
\newcommand{\dd}{\,\mathrm{d}}
\begin{document}
\section*{UnsolvedMath ID: 800067}
\textbf{Source problem code:} catalogue entry
\subsection{Definitions and mathematical statement}
Let \(M\) be a closed hyperbolic \(3\)-manifold with \(b_1(M):=\dim_{\R}H^1(M;\R)\geq 1\). The formulation below uses the usual oriented-manifold and cooriented-contact-structure interpretation. The supplied record does not explicitly impose these conventions, so their applicability requires review.

Under this interpretation, fix an orientation of \(M\). For an integral homology class \(z\in H_2(M;\Z)\), the Thurston norm is
\[
x_M(z)=\min_{[S]=z}\sum_i\max\{0,-\chi(S_i)\},
\]
where \(S=\bigsqcup_i S_i\subset M\) ranges over embedded closed oriented surfaces representing \(z\). Extend this norm to \(H_2(M;\R)\) by homogeneity and continuity. For \(a\in H^2(M;\R)\), define the dual Thurston norm by
\[
x_M^*(a)=\sup\bigl\{|\langle a,z\rangle|:z\in H_2(M;\R),\ x_M(z)\leq 1\bigr\}.
\]
Let \(r:H^2(M;\Z)\to H^2(M;\R)\) be change of coefficients. Interpret an \emph{even integral real cohomology class} as a class in \(2r(H^2(M;\Z))\). This interpretation does not select an integral lift or its torsion component.

A cooriented contact structure is a plane distribution \(\xi=\ker\alpha\subset TM\), where \(\alpha\) is a global smooth \(1\)-form satisfying \(\alpha\wedge\mathrm d\alpha\neq 0\) everywhere. The orientation of \(M\) and the coorientation of \(\xi\) orient \(\xi\), giving its Euler class \(e(\xi)\in H^2(M;\Z)\). A contact structure is \emph{tight} if it contains no overtwisted disk.

The supplied question, under these conventions, asks whether for every such \(M\) and every \(a\in H^2(M;\R)\),
\[
\bigl(a\in 2r(H^2(M;\Z))\ \text{and}\ x_M^*(a)=1\bigr)
\quad\Longrightarrow\quad
\exists\,\xi\ \text{tight and cooriented on }M
\quad r(e(\xi))=a.
\]
The equality here is in real cohomology. The source does not specify whether a particular integral lift must be realized, whether contact structures must be positive relative to a chosen orientation, or whether orientation and coorientation are implicit hypotheses. These qualifications remain unresolved; the display records a conventional reading rather than a verified clarification of the source.
\subsection{Short English statement}
For every closed hyperbolic \(3\)-manifold with positive first Betti number, must every even integral real cohomology class on the boundary of the dual Thurston norm unit ball be the real Euler class of a tight contact structure? Orientation, coorientation, positivity, and the meaning of integral realization require clarification from the source.
\subsection{Sources}
\begin{itemize}
\item Local UnsolvedMath registry: permanent ID \(800067\), exact source problem code \texttt{catalogue entry}, title ``Foliations and Diffeomorphism Groups --- Question 8.'' Supplied upstream revision: \texttt{d371ee917ac9647fbb06cd6fc529deec6f3dda82}.
\item Nariman and Yazdi, \emph{Problem List on Foliations and Diffeomorphism Groups} (2025), Question 8. Supplied source URL: \url{https://arxiv.org/abs/2503.18274}. Attribution and year are as supplied in the registry.
\item The record reports extraction by \texttt{arxiv\_latex} and access on 2026-10-16. This access date is supplied provenance, not an independently verified source visit. The registry labels the problem ``open''; that label is supplied curation, not an independently verified literature assessment.
\end{itemize}
\end{document}
